\section{Liquidity Protocol Integration}
\label{sec:liquidity-integration}

The \Liquidity{} Protocol is the compliance-by-construction settlement
standard for tokenised securities formally launched by \Liquidity{}.io
on April 1, 2026 with a soundness proof published alongside
\cite{liquidity-protocol-formal}. \Lux{}, \Hanzo{}, \Zoo{}, and Pars
adopted the \Liquidity{} Protocol on April 20, 2026; this section
documents \Zoo{}'s adoption \cite{zoo-adopts-liquidity-protocol}.

\subsection{What the Liquidity Protocol provides}

\Liquidity{} Protocol is the contract a chain accepts to integrate
with the \Liquidity{}.io control plane. The contract specifies:

\begin{enumerate}[leftmargin=1.4em]
  \item The KYC/AML perimeter: where checks happen, what they verify,
        what evidence is admissible.
  \item The broker-dealer interface: which actions require BD
        signoff, what the BD's signing key looks like on-chain, how
        BD attestations flow through transactions.
  \item The transfer-agent interface: how the official record is
        updated, what events constitute a record update, how the
        chain's view of beneficial ownership is reconciled with the
        TA's view.
  \item The market-center interface: what data the venue must
        publish (best-bid-best-offer, last-trade, volume), what
        venue-side reporting goes back to the BD and TA.
  \item The dispute interface: how a flagged transaction is unwound,
        what cancel/reverse semantics look like at the chain level.
\end{enumerate}

The soundness proof in \cite{liquidity-protocol-formal} establishes
that an integrating chain that follows the contract cannot construct
a state that is unattributable, unrecoverable, or non-compliant
under SEC ATS, FINRA BD, and registered TA rules. Zoo inherits this
soundness by integrating against the formal contract.

\subsection{The Liquidity precompile on Zoo D-Chain}

\Zoo{} D-Chain ships a precompile at a reserved address that
implements the chain-side half of the \Liquidity{} Protocol. The
precompile interface exposes:

\begin{lstlisting}[language=C++,caption={Liquidity precompile (Zoo D-Chain)}]
namespace zoo::dchain::liquidity {

  // KYC/AML perimeter
  bool is_kyc_clear(Address holder);
  bool is_aml_clear(Address holder, Hash transaction_intent);

  // Broker-dealer attestations
  bool bd_authorised(Address holder, ActionId action);
  Hash bd_signoff_required_for(ActionId action);

  // Transfer agent
  void notify_ta_of_trade(TradeEvent ev);
  Hash ta_record_root_at_block(BlockHeight h);

  // Market center
  void publish_bbo(Pair p, Price bid, Price ask);
  void publish_last_trade(Pair p, Price px, Quantity qty);

  // Dispute / unwind
  bool is_action_disputed(Hash action_hash);
  void execute_unwind(Hash action_hash, UnwindReason reason);
}
\end{lstlisting}

The precompile is the only on-chain entity that can mutate the
\texttt{SecurityClass} state for a regulated security. The matching
engine, the AMM, the order types --- all of them call into the
precompile to verify each action before applying state changes.
Trade requests for which the precompile returns ``not allowed''
never reach the order book.

\subsection{Cross-chain trade routing through B-Chain}

For pairs where the deepest liquidity lives outside \Zoo{} D-Chain
(e.g., a security listed on \Zoo{} but with cross-listing liquidity
on \Lux{} C-Chain or external venues), trades route via B-Chain
\cite{lp-134}:

\begin{enumerate}[leftmargin=1.4em]
  \item Order submitted on \Zoo{} D-Chain.
  \item Matching engine determines best route includes a cross-chain
        leg.
  \item B-Chain bridge message emitted; \Lux{} \DEX{} \cite{lux-dex}
        executes the cross-chain leg.
  \item Settlement leg returns via B-Chain; the cross-chain attestation
        carries triple-cert proof.
  \item Local fill is recorded; transfer-agent notify fires; trade
        finalised.
\end{enumerate}

The \Liquidity{} precompile is consulted on both sides of the
B-Chain hop: \Zoo{} D-Chain enforces \Zoo{}-side compliance; \Lux{}
C-Chain (also a \Liquidity{} Protocol integrator) enforces \Lux{}-side
compliance. A trade that violates compliance on either side is
rejected.

\subsection{Why a precompile and not a smart contract}

The \Liquidity{} integration must be unbypassable. If KYC checks
lived in an EVM contract, an application contract could fork the
precompile, omit the KYC call, and trade against the matching engine
directly. By living in the precompile (chain-built-in code, not
EVM-deployable), the \Liquidity{} integration is enforced at the
chain interpreter level. There is no application-contract path that
the matching engine accepts that does not first go through the
precompile.

This is the compliance-by-construction property that the \Liquidity{}
Protocol soundness proof relies on.

\subsection{Soundness inheritance}

The \Liquidity{} Protocol formal proof \cite{liquidity-protocol-formal}
shows that a chain that:

\begin{enumerate}[leftmargin=1.4em]
  \item exposes the \Liquidity{} precompile as the only mutator of
        \texttt{SecurityClass} state, and
  \item routes every trade, transfer, and corporate action through
        the precompile, and
  \item maintains the cross-chain invariants for B-Chain hops,
\end{enumerate}

inherits the soundness property: no chain state is reachable that
violates the SEC ATS, FINRA BD, or registered TA obligations
encoded in the protocol. Zoo D-Chain meets all three preconditions
and therefore inherits soundness from April 20, 2026 onward.

\subsection{Operational interfaces}

\begin{itemize}[leftmargin=1.2em]
  \item \textbf{Issuer onboarding.} Issuer registers with
        \Liquidity{}.io; the resulting attestation is published on
        A-Chain; \Zoo{} D-Chain's listing-admission flow reads the
        A-Chain attestation; once verified, the listing becomes
        active.
  \item \textbf{Holder onboarding.} Holder onboards via wallet UX
        through \Liquidity{}.io's KYC flow; the resulting on-chain
        record sets a holder-eligibility bit on A-Chain that the
        precompile reads.
  \item \textbf{Audit interface.} Regulators query the on-chain
        record directly; the chain log is the audit trail; no
        secondary database is needed.
\end{itemize}
